% Related Work section expanded for the moderate-length revision.
% This file is a section fragment, not a standalone LaTeX document.

\section{Related Work}
\label{sec:related_work}

\subsection{Uniform Quantization and Integer-Only Inference}
Uniform affine quantization maps floating-point tensors to a finite set of integer values using a scale and, when needed, a zero-point. This formulation is widely used because it gives a simple connection between floating-point tensors and integer arithmetic. Integer-only inference further requires the computational graph to avoid repeated floating-point operations during deployment. Jacob \textit{et al.} showed how quantized neural networks can be executed with integer-arithmetic-only inference \cite{jacob2018quantization}, while later work such as HAWQ-V3 further studied integer-only and mixed-precision quantization under hardware-related constraints \cite{yao2020hawqv3}. This paper follows the uniform quantization setting, but evaluates it through fake quantization rather than a real integer inference kernel.

\subsection{Post-Training Quantization for CNNs}
PTQ determines quantization parameters after training and usually relies on a calibration subset. General surveys and white papers summarize observer design, calibration, quantization granularity, and reconstruction-based evaluation for neural network quantization \cite{gholami2021survey,nagel2021whitepaper,krishnamoorthi2018quantizing}. Several PTQ methods improve over direct rounding or simple calibration by optimizing local reconstruction objectives. LAPQ studies loss-aware layer-wise parameter selection \cite{nahshan2019loss}, AdaRound adapts weight rounding decisions using a local loss approximation \cite{nagel2020adaround}, and BRECQ reconstructs network blocks to improve low-bit PTQ behavior \cite{li2021brecq}. ZeroQ addresses the related case where original calibration data are unavailable by synthesizing data from batch-normalization statistics \cite{cai2020zeroq}. The present work is more limited: it does not optimize weights or blocks, and it focuses only on activation threshold selection in a simulated PTQ pipeline.

\subsection{Low-Bit and 4-Bit Quantization}
Reducing precision from 8 bits to 4 bits increases the impact of both rounding and clipping error. Early low-bit PTQ work showed that 4-bit quantization is feasible but more fragile than 8-bit quantization, especially when quantization parameters are selected without retraining \cite{banner2018post,choukroun2019low}. Later methods such as AdaRound and BRECQ improve 4-bit PTQ by changing how weights or blocks are reconstructed \cite{nagel2020adaround,li2021brecq}. HAWQ-V3 also studies mixed-precision choices and integer-only deployment constraints \cite{yao2020hawqv3}. These works provide context for the difficulty of INT4 PTQ. In contrast, this paper does not attempt to build a comprehensive low-bit benchmark or claim state-of-the-art performance; it examines a simple activation clipping rule under fixed experimental conditions.

\subsection{Activation Clipping and Outlier Suppression}
Activation quantization is often sensitive to outliers because a small number of large values can determine the full quantization range under MinMax calibration. Parameterized activation clipping learns clipping thresholds during training \cite{choi2018pact}, while outlier channel splitting reduces the effect of outlier channels without retraining \cite{zhao2019outlier}. SmoothQuant, although developed for large language models rather than CNNs, is also relevant because it explicitly treats activation outliers as a source of quantization difficulty and migrates part of that difficulty from activations to weights \cite{xiao2022smoothquant}. The method in this paper follows the broader idea of reducing outlier sensitivity, but it uses no retraining, no learnable clipping parameters, and no equivalent weight transformation. Each post-ReLU activation threshold is selected from a predefined percentile set according to calibration reconstruction MSE.
